\section{Evaluating Pillar 1: Extensibility}
\label{sec:foundational-ingestion-extensibility}

This section addresses \textbf{RQ1}: whether JD-driven, recruiter-configurable metrics 
and reliable ingestion can support downstream scoring.
In order for a screening system to be effective, it requires data to be ingested in a way that is accurate and reliable.
Therefore, the first area that we will evaluate is MERIT's ingestion layer, to ensure that it is reliable enough to
parse unstructured documents, specifically CVs and JDs. We do not consider structured data sources such as LinkedIn and GitHub
since there is no need to parse them, as they are already in a structured format. We begin with two studies, 
one for CVs and one for JDs, to evaluate the accuracy of the ingestion layer. Both studies compare the system's extraction 
outputs against human-verified ground-truth datasets, evaluating a total of 30 CVs (balanced across diverse layouts) and 15 JDs.
To finish the evaluation of this section, we then consider how well the metric priority predictions align with that
of a panel of recruiters, testing whether TF-IDF-derived priorities avoid the fixed-coefficient rigidity criticised in prior work 
by Gugnani and Misra~\cite{gugnani2020} and Lo et al.~\cite{lo2025aihiring}. Graphical representations for this section are 
in Appendix~\ref{appendix:ingestion-charts}.

\subsection{Evaluation Metrics}
\label{subsec:mathematical-metrics}

For our evaluation, we will be evaluating both set-based and string-based comparisons between the ground truth,
and the extracted data for our three core categories (skills, experience and identity).

\subsubsection{Skill Set Metrics}
Given $S_E$ as the set of extracted skills and $S_{GT}$ as the set of ground truth skills, we can
formulate the precision ($P_{skills}$), recall ($R_{skills}$), and F1-score ($F1_{skills}$) as follows:
{\small
\begin{align}
    P_{skills}  &= \frac{|S_E \cap S_{GT}|}{|S_E|}, &
    R_{skills}  &= \frac{|S_E \cap S_{GT}|}{|S_{GT}|}, &
    F1_{skills} &= \frac{2 P_{skills} R_{skills}}{P_{skills} + R_{skills}}
\end{align}
}

\subsubsection{Experience Recall}
For experience extraction, we mainly care about recall, since if we miss an experience, we might miss out on a
correct match. False positives are not a significant concern for this metric due to the strict nature of experience
extraction.

Given a set of ground truth experiences $E_{GT}$, and a set of extracted experiences $E_E$, we define a
successful match as follows:
\begin{equation}
    \text{Match}(e_{gt}) = \mathbb{I}\left(\exists e_e \in E_E \text{ s.t. } \text{norm}(e_{gt}) \subseteq \text{norm}(e_e)\right)
\end{equation}

Experience Recall ($R_{exp}$) is thus:
\begin{equation}
    R_{exp} = \frac{1}{|E_{GT}|} \sum_{e_{gt} \in E_{GT}} \text{Match}(e_{gt})
\end{equation}

\subsubsection{Identity and Metadata Matching}
When it comes to the extraction of identity for both CVs and JDs, we only consider the exact matches, since
metrics such as precision and recall are not applicable in that case. For CVs we do this for both name and 
email addresses. For the JD, while we do perform a similar operation, we are less strict with the matching,
it is considered a match if the ground truth label is a substring of the extracted label to account for the 
difficulty of matching the company name exactly in JDs.

\subsection{Study 04: CV Ingestion Accuracy}
\label{subsec:study-04}

We first evaluate the accuracy of the CV parser, extracting candidate data from 15 single-column CVs. We also
replicate this data across 15 multi-column CVs, ensuring that the content is the same for a given candidate,
to also evaluate the impact of a CV's layout on the accuracy of the parser. The CVs are designed to be diverse
in nature, across different roles and seniority levels. We craft a ground truth dataset for each candidate,
with the author verifying the accuracy and completeness of these JSON payloads.
Table~\ref{tab:cv_parser_results} presents the compiled results.

\begin{table}[htbp]
    \centering
    \caption[Study~04: CV ingestion by layout]{Study~04: CV ingestion accuracy by layout ($N=30$).}
    \label{tab:cv_parser_results}
    \merittablesize
    \setlength{\tabcolsep}{6pt}
    \begin{tabular}{l|c|c|c}
        \hline
        \textbf{Metric} & \textbf{Single-col.\ ($N{=}15$)} & \textbf{Multi-col.\ ($N{=}15$)} & \textbf{Overall ($N{=}30$)} \\
        \hline
        Skill Extraction Precision   & 69.6\%  & 69.6\%  & 69.6\%  \\
        Skill Extraction Recall      & 38.6\%  & 38.6\%  & 38.6\%  \\
        Skill Extraction $F_1$-Score & 49.4\%  & 49.4\%  & 49.4\%  \\
        Experience Recall            & 34.3\%  & 44.7\%  & 39.5\%  \\
        Name Match Rate              & 53.3\%  & 13.3\%  & 33.3\%  \\
        Email Match Rate             & 100.0\% & 100.0\% & 100.0\% \\
        \hline
    \end{tabular}
\end{table}

The distribution of individual scores across the template pool is plotted in
Figure~\ref{fig:ingestion-combined} (panel~a, Appendix~\ref{appendix:ingestion-charts}).

\subsubsection{Qualitative Analysis}

We have four key findings from this study, each with different root causes.

The first finding is that skill extraction is not hindered by the layout of a CV, but rather the dictionary that it is 
using to identify the skills. We can observe this as both layouts achieve the same precision, recall and F1-score ($69.6\%$, $38.6\%$ and $49.4\%$ respectively).
The source that we collect skills from, Devicon, catalogues programming languages as well as frameworks, however, it completely
ignores categories that are common on real engineering CVs. One such example is tooling for observability, such as Prometheus and Grafana.
The precision indicates that what is found is correct for the most part, but the limiting issue is the dictionary we use, rather than the 
logic of our extraction process. This finding is consistent with resume-parsing benchmarks such as ResumeBench \cite{emnlp2025resumebench}.

Another important observation is that the detection of candidate name drops sharply under multi-column layouts.
There is a 40\% drop between single-column and multi-column layouts when it comes to identifying the candidate name.
In a typical single-column layout, candidate name is generally the first item in the CV, at the top of the document.
However, for multi-column layouts, this is often not the case. This is a clear example of structural failure in the 
ingestion layer due to our parser having simplistic heuristics for name extraction.

We also find that email extraction is fully robust, achieving a 100\% match rate across both layout types,
confirming that our regex-based extraction process is not hindered by the layout of the candidate CV.

Finally, we end with counterintuitive results for experience extraction, noting that experience recall rises from
34.3\% in single-column layouts to 44.7\% in multi-column layouts. The structural difference between the two layouts
is definitely the reason, however, initially it was not clear why it was higher in the multi-column layout. Further
inspection leads us to believe that the reason is due to Single-Column layouts presenting experience as inline, with
multiple pieces of information on the same line, separated by a pipe character, or other delimiters. This must
be parsed as a single unit before we can isolate details such as the employer name and the role. Multi-column
layouts tend to have less width due to the introduction of more columns, so candidates tend to have these details on
separate lines. This actually makes it easier for the parser to isolate details, which could suggest why multi-column
layouts have better recall for experience compared to single-column layouts.

CVs are unstructured documents, and as such, there is a need for structured, schema-based data sources to be able
to account for parsing errors that may arise from CV ingestion (such as a LinkedIn profile). In real-world
applications, it is likely that commercial Applicant Tracking Systems (ATS) suffer from the same structural
parsing failures that we have observed in our study. The black-box nature of these screening systems would 
make it difficult for recruiters to know if this is actually the case for why a candidate receives a low score.
It is for this reason that candidates are strongly advised to use simple, single-column templates, as discussed in practitioner ATS guidance \cite{marr2019}.
Recruitment-management systems also tend to discard applicants who fail to mirror job-description keywords and criteria \cite{fuller2021hiddenworkers}.
In turn this uncovers a systemic bias: automated screening systems penalise highly qualified candidates if the
layout of their CV confuses the parser, consistent with the parsing failures we observed.

\textbf{Outcome:} Does not meet the internal parsing $\geq 80\%$ target (skill $F_1 = 49.4\%$, name match = $33.3\%$).
Precision is acceptable, but dictionary recall and multi-column name heuristics bound RQ1 ingestion.

\subsection{Study 05: Job Description Accuracy}
\label{subsec:study-05}
The JD parser was evaluated in a similar manner, with a total of 15 job descriptions across different software
engineering roles to ensure a wide range of formats and technical requirements. The compiled mean results
are summarised in Table~\ref{tab:jd_parser_results}, the per-document distribution is in
Figure~\ref{fig:ingestion-combined} (panel~b, Appendix~\ref{appendix:ingestion-charts}).

\begin{table}[htbp]
    \centering
    \caption[Study~05: JD extraction accuracy]{Study~05: mean JD extraction accuracy ($N=15$).}
    \label{tab:jd_parser_results}
    \merittablesize
    \setlength{\tabcolsep}{6pt}
    \begin{tabular}{l|l|r}
        \hline
        \textbf{Metric} & \textbf{Evaluation focus} & \textbf{Mean (\%)} \\
        \hline
        Skill Requirements Precision   & Target skill relevance       & 60.02\% \\
        Skill Requirements Recall      & Target skill coverage        & 34.13\% \\
        Skill Requirements $F_1$-Score & Overall requirement fidelity & 40.81\% \\
        Job Title Match Rate           & Role identification        & 46.67\% \\
        Company Match Rate             & Employer identification    & 6.67\%  \\
        \hline
    \end{tabular}
\end{table}

\subsubsection{Qualitative Analysis \& Strategic Ingestion Failure Cases}

Unlike CVs, which serve as concise professional summaries, Job Descriptions (JDs) are inherently
marketing-oriented documents designed to attract candidates. As a result, they are often dense with corporate
boilerplate, such as institutional missions, values, wellness benefits, and equal opportunity clauses.
This structural noise significantly degrades parsing accuracy. Because section headings are frequently
styled with generic tags (e.g.\ ``About Us'' or ``Who We Are''), the parser struggles to isolate core technical
qualifications from general corporate copy. This lexical contamination is directly reflected in the metrics,
driving a low Skill Recall of 34.13\% and restricting Skill Precision to 60.02\%, as the extraction engine
fails to separate critical technology requirements from marketing jargon.

The most notable observation within this evaluation is the poor performance of company name extraction, with a match rate
of 6.67\%, only successfully extracting one of the fifteen JD company names. The reason for this is structural, rather than 
semantics. Employers often represent company names in text that is not easily extractable by a parser, for example, they
represent their name as a logo or image. Text extractors (in our case, pdfminer) cannot extract this textual information from
an image, it is therefore forced to rely on the ``About Us'' section. If the name is not present in this section, or 
referred to via generic pronouns (such as ``Our company''). Then the extraction process fails. Fortunately, this is a non-critical
bottleneck, since the only usage of this metric is for the recruiter to know the company name, and not for the scoring process.
It is a convenience feature rather than a core functional requirement.

\textbf{Outcome:} Does not meet parsing $\geq 80\%$ (skill $F_1 = 40.81\%$).
While employer extraction is not critical for ranking, the JD skill recall still fails to meet RQ1's success criteria.

\subsection{Study 13: Metric Prediction vs Recruiter Alignment}
\label{subsec:study-13-metric-prediction}

While the past two studies evaluate MERIT's extraction, this study instead evaluates the weighting side of RQ1. We
aim to test whether TF-IDF derived metric priorities are able to align with expert consensus, as opposed to merely
having a fixed coefficient scheme.
We test \textbf{H1} (positive Spearman $\rho$ within each role) and \textbf{H2} (low MAE versus consensus on MERIT's 1--5 scale,
where 1 = essential and 5 = optional).

\subsubsection{Method}
When generating the market corpus, we employ Tonic.ai to synthesise a dataset of 100 JDs, providing a relatively realistic
dataset to aid in this evaluation. We have three hold-outs which are excluded from the generation of the corpus, that being a Junior 
Frontend, senior ML and Lead DevOps role. This is to prevent TF-IDF from being biased towards these roles. We have one global corpus,
we do not partition by seniority nor by role family.

On the testing side, we have two recruiters and two academic peers who are tasked with weighting the metrics given a JD. We intentionally
hide the suggestions that TF-IDF would have suggested to ensure fairness. Our consensus is determined by the median of these ratings, and 
we compare these results to TF-IDF suggestions, calculating the exact match, mean absolute error (MAE) and Spearman $\rho$.
The graphical representations for this study are available in Appendix~\ref{appendix:metric-prediction-charts}.

\subsubsection{Results}
As shown in Table~\ref{tab:study-13-metric-alignment}, TF-IDF is not well
calibrated for exact matches. We can see that none of the roles manage $\geq 50\%$ for the exact match rate between 
TF-IDF and recruiter consensus. However, we can observe that in most cases, Spearman $\rho$ is positive, showing some
alignment between the TF-IDF predictions and the recruiter consensus. This is a good sign, as it 
indicates that the TF-IDF predictions are generally in the right direction, but not necessarily a perfect match.
\begin{table}[htbp]
    \centering
    \caption[Study~13: TF-IDF vs consensus]{Study~13: TF-IDF vs.\ consensus ($N=4$, $n=27$ skill--role pairs).}
    \label{tab:study-13-metric-alignment}
    \merittablesize
    \setlength{\tabcolsep}{6pt}
    \begin{tabular}{l|c|c|c|c}
        \hline
        \textbf{Hold-out role} & \textbf{Skills} & \textbf{Exact match} & \textbf{MAE} & \textbf{Spearman $\rho$} \\
        \hline
        Junior Frontend (React) & 9  & 44.4\% & 0.78 & 0.52 \\
        Senior ML Engineer       & 8  & 25.0\% & 1.88 & $-0.45$ \\
        Lead DevOps / Platform   & 10 & 20.0\% & 1.20 & 0.40 \\
        \hline
        \textbf{Aggregate} ($n{=}27$) & 27 & 29.6\% & 1.26 & 0.25 \\
        \hline
    \end{tabular}
\end{table}

Of the 27 skill--role pairs, eight achieve exact agreement and a further nine differ by only one Likert band (63\% within $\pm 1$).
The near misses show that our metric predictions should be treated as reasonable guidance, rather than simply a default value.


\subsubsection{Limitations and their Implications}
\label{subsubsec:study-13-limitations}

This is a pilot study ($N = 4$ evaluators, $n = 27$ pairs), and the results are not generalisable to the wider market.
We are not able to conduct a larger panel study due to the time and resource constraints of this project.

Another limitation is due to the way we create the corpus, or rather, the fact that we only use one corpus for all TF-IDF calculations.
Due to this, we encounter cross-role IDF interference. To take an example, let's consider Python, a language used in many different roles such
as backend development, machine learning, data science and more. Since Python is common in these roles, it will inevitably be in a significant
proportion of job descriptions and, as a result, we would have a high IDF score, regardless of the role, making it seem like Python is 
always a common skill. This down-ranks Python on the metric predictions, regardless of the role.

In this study, we also did not ask our participants to spread their ratings evenly between the 1 and 5 scale. Compare this to TF-IDF
which has an in-built min--max normalisation to the range $[1.0, 5.0]$ which ensures a relatively even distribution of these ratings.
This disparity in rating distribution is likely to play a role in our evaluation results, particularly the exact match rate. 
That mismatch can depress exact-match rate and $\rho$ without invalidating the pilot, it is a
comparison caveat we report alongside the within-$\pm 1$ band counts above. Cross-study validity threats are summarised in
Section~\ref{sec:threats-to-validity}.

\textbf{Outcome:} Exploratory RQ1 criterion only: aggregate Spearman $\rho = 0.25$ with mixed hold-outs 
(Table~\ref{tab:study-13-metric-alignment}).
\textbf{H1} is partly supported and \textbf{H2} is not met. TF-IDF priorities should be treated as guidance, 
not automatic weights, with recruiter overrides on the frontend.

\subsection{Downstream Mitigation}
\label{subsec:downstream-mitigation}

These quantitative findings empirically justify two of MERIT's three core architectural pillars.

\subsubsection{Multi-Source Fusion (Pillar 2)}

The recall of 38.6\% means that CV-only screening will often miss the majority of a candidate's technical competencies
under the Devicon dictionary. The drawbacks of a dictionary-based approach are clearly visible here, as well as the 
reliance on a single source of truth for areas such as skill extraction. By ingesting other sources into candidate
scoring (such as GitHub repository language statistics and LinkedIn skill endorsements), MERIT can recover 
skills that the CV parser may have failed to recognise. 


\subsubsection{Glass-Box Explainability (Pillar 3)}
The issues presented above would often fail to be surfaced to recruiters in typical black-box scoring systems.
This highlights the need for a system that can surface these parsing failures to recruiters. The addition of 
MERIT's explainability layer achieves this by providing in-depth reviews of both the candidate's original 
data, as well as what MERIT was able to extract from it. This ensures that recruiters are able to audit and alter
their decisions based on the information that MERIT was able to gather. For example, a recruiter may observe that
a metric, such as Python, may not have been extracted from the candidate CV due to a parser error. The recruiter would
then be able to either observe the original CV to challenge the claim, or perhaps consider other data sources.